\section{Full scan}
\label{sec:scan}
Approach A is the reference candidate-search method. Once clusters and eligible rows are chosen,
it scores every eligible row. There is no approximation \emph{inside} that selected pool.
The loop repeats the same operations and can read packed codes in order when filtering does
not leave scattered rows.

\fig{scan}{Approach A. The selected pool is smaller than the corpus only because of routing or filters.}{1}

\Needspace{12\baselineskip}
\subsection{Data layout}
The common index from Equation~\eqref{eq:base} is sufficient. Cluster $j$ occupies the row
range between offsets $j$ and $j+1$. A row contains its packed code, and the corresponding ID
array maps the row back to the original document.

\par\noindent\begin{minipage}{\linewidth}
\begin{lstlisting}[caption={Candidate search for Approach A.}]
ScanSelectedClusters(q, selected_clusters, eligible, table, R):
    best = empty top-R heap
    for cluster j in selected_clusters:
        for row in eligible rows of j:
            score = ScoreDocument(code[row], table)
            best.offer(score, document_id[row])
    return best in descending score, breaking ties by ID
\end{lstlisting}
\end{minipage}\par

The loop is over $\Nelig$ rows, not necessarily $NP/C$. If all clusters are searched and no
filter removes documents, $\Nelig=N$. Empty filtered clusters require no document scoring,
although finding that they are empty can still have a cost.

\Needspace{12\baselineskip}
\subsection{Work count}
\begin{center}\small
\begin{tabularx}{\linewidth}{@{}P{.25\linewidth}YP{.23\linewidth}@{}}\toprule
Step & Work & Live temporary data\\\midrule
Locate cluster ranges & $P$ offset/range accesses & selected cluster IDs\\
Visit eligible rows & $\Nelig$ row iterations & current row/iterator\\
Score rows & $\Nelig G$ lookup terms & shared score table, accumulator\\
Keep best $R$ & $H_s(\Nelig,R)$ selection operations & top-$R$ heap\\
Return candidates & sort at most $R$ results if needed & result array or heap storage\\\bottomrule
\end{tabularx}\normalsize\end{center}

Let $c_{\rm range}$ denote range setup and $T_{\rm rank}(R)$ denote final candidate ordering. Then
\begin{equation}
T_{A,\rm candidates}=P c_{\rm range}+\Nelig G c_{\rm lut}
                    +H_s(\Nelig,R)c_{\rm heap}+T_{\rm rank}(R).
\label{eq:scan-time}
\end{equation}
The general end-to-end latency is
\begin{equation}
\boxed{T_A=\Tcommon+T_{A,\rm candidates}+\Tfinish(R').}
\label{eq:scan-total}
\end{equation}
For the stated heap implementation, the candidate-stage time complexity is
\[
O\!\left(P+\Nelig\ceil{d/g}+\Nelig\log(R+1)+R\log(R+1)\right).
\]
The lookup term count is not a claim that each term is exactly one CPU instruction.

\Needspace{12\baselineskip}
\subsection{Data reads and memory}
The number of packed-code data bytes read for scoring is
\begin{equation}
V_{A,\rm codes}=\Nelig\,s_{\rm code}.
\end{equation}
ID reads, filter data, table accesses and result-heap accesses are additional traffic. A cache
may satisfy many of them. This formula therefore counts only code data; it does not measure all reads from main memory.
Sequential and scattered rows may have different effective $c_{\rm lut}$ or $c_{\rm row}$.

Index storage and temporary memory during candidate search are
\begin{align}
M_{A,\rm index}&=\Mbase,\\
Q_{A,\rm search}&=\Qbase+Q_{\rm row\ iterator}+Q_{\rm scan\ workspace}.
\label{eq:scan-memory}
\end{align}
A streaming scan does not need an $N$-element score array. A matrix-based batch implementation
may choose to allocate such an array; then that allocation must be added. Two implementations can therefore return the same ranking while using different amounts of
peak RAM.

\Needspace{12\baselineskip}
\subsection{Routing example}
The eight documents from Section 1 are partitioned into
\[
\mathcal C_0=\{1,3,4,6\},\qquad \mathcal C_1=\{2,5,7,8\}.
\]
This is an illustrative fixed partition, not a claim about the output of a particular
clustering run. Suppose the router opens only $\mathcal C_1$. A computes scores
$1.2,0.6,0.8,-0.2$ and returns documents 2 and 7. The global top two are 1 and 2, so recall is
$1/2$. The scan is exact within the opened cluster; document 1 is excluded by routing.

Opening both clusters returns documents 1 and 2. A scan can therefore be exact within its selected
pool while its global recall remains below one.

\Needspace{12\baselineskip}
\subsection{SIMD}
SIMD means one instruction operates on several values at once. It can accelerate code extraction,
lookup accumulation, and selection. It does not by itself change the number of stored bits.
There are two useful ways to model its effect.

\textbf{Calibrated model.} The scoring routine's time per work unit,
$c_{\rm lut}$, is measured for the relevant dimensions and access pattern. If that rate already includes
selection, a combined $T_{\rm score+select}$ excludes a separate heap-cost term to avoid double counting.

\textbf{Resource model.} If a scoring routine performs $U_{\rm op}$ compute units and causes $V_{\rm DRAM}$ bytes of DRAM
traffic, compute throughput $\Gamma$ and memory bandwidth $\beta$ give the lower bound
\begin{equation}
T_{\rm kernel}\geq\max\{U_{\rm op}/\Gamma,\ V_{\rm DRAM}/\beta\}.
\label{eq:resource}
\end{equation}
An effective SIMD factor $s_{\rm simd}$ can change $\Gamma$ to $s_{\rm simd}\Gamma_0$.
It cannot make the bandwidth term disappear. Random access, synchronization and fixed overhead
can put actual time above the bound. This resource estimate is not added to a fitted
coefficient that already includes the same work.

For multiple queries, a float score matrix can be written as $QB^\top$, where rows of $B$ are
expanded binary codes. Matrix operations can compute this product, but code expansion, batch formation, device
transfers and score-buffer storage incur costs. Fast GPU batch processing does not necessarily imply low
single-query latency.

\begin{resultbox}{Accuracy}
For a fixed eligible pool, A returns its exact binary top $R$ under the stated tie rule.
With every eligible cluster searched, it has 100\% binary recall. With partial routing,
its global binary recall is limited by which true neighbors reach the selected pool.
The scan cannot restore documents missed by routing, quantization or an earlier filter.
\end{resultbox}
